% chemfrag-doc.tex -- user manual of the chemfrag package
% Copyright (C) 2026 Vaibhav Blayer -- LPPL 1.3c or later
\documentclass[11pt]{article}
\usepackage[T1]{fontenc}
\usepackage{lmodern}
\usepackage[a4paper,margin=2.4cm,headheight=14pt]{geometry}
\usepackage{microtype}
\usepackage{xcolor}
\usepackage{chemfrag}
\usepackage{booktabs,array,longtable}
\usepackage[most]{tcolorbox}
\usepackage{enumitem}
\usepackage{fancyhdr}
\usepackage[hidelinks]{hyperref}

% ---------------------------------------------------------------- colours
\definecolor{cfaccent}{HTML}{1F6F8B}
\definecolor{cfback}{HTML}{F4F8FA}
\definecolor{cfrule}{HTML}{C9DCE3}
\definecolor{cfcode}{HTML}{8A3B12}
\definecolor{cfgray}{HTML}{6B7780}
\hypersetup{colorlinks,linkcolor=cfaccent,urlcolor=cfaccent,citecolor=cfaccent}

% ---------------------------------------------------------------- layout
\pagestyle{fancy}
\fancyhf{}
\fancyhead[L]{\small\color{cfgray}\textsf{chemfrag} v\chemfragversion}
\fancyhead[R]{\small\color{cfgray}\nouppercase{\leftmark}}
\fancyfoot[C]{\small\color{cfgray}\thepage}
\renewcommand\headrulewidth{0pt}
\setlength\parindent{0pt}
\setlength\parskip{0.55em}
\setlist{itemsep=0.2em,topsep=0.3em}

\usepackage{titlesec}
\titleformat{\section}{\Large\sffamily\bfseries\color{cfaccent}}{\thesection}{0.7em}{}
\titleformat{\subsection}{\large\sffamily\bfseries}{\thesubsection}{0.6em}{}
\titlespacing*{\section}{0pt}{1.6em}{0.6em}

% ---------------------------------------------------------------- code
\lstdefinestyle{cf}{
  basicstyle=\ttfamily\small,
  columns=fullflexible,
  keepspaces=true,
  breaklines=true,
  upquote=true,
  moredelim=[is][\color{cfcode}]{|}{|},
}
\tcbset{
  cfbox/.style={
    enhanced, colback=cfback, colframe=cfrule, boxrule=0.5pt, arc=2pt,
    left=6pt, right=6pt, top=4pt, bottom=4pt,
    listing options={style=cf},
  },
}
% code on the left, result on the right
\newtcblisting{example}{
  cfbox, listing side text, righthand width=0.42\linewidth,
  halign lower=center, valign lower=center,
  segmentation style={draw=cfrule,solid}
}
% code above, result below (wide molecules)
\newtcblisting{examplewide}{
  cfbox, text above listing=false, listing above text,
  halign lower=center, segmentation style={draw=cfrule,solid}
}
% code only
\newtcblisting{code}{cfbox, listing only}

\newcommand\cs[1]{\texttt{\textbackslash#1}}
\newcommand\meta[1]{\ensuremath\langle\textit{#1}\ensuremath\rangle}
\newcommand\key[1]{\texttt{\color{cfcode}#1}}
\newcommand\frag[1]{\texttt{!\{#1\}}}
\newcommand\pkg[1]{\textsf{#1}}

% ring gallery
\newcommand\ringnum[1]{\scriptstyle\color{cfaccent}#1}
\newcommand\ringcell[2]{%
  \begin{minipage}[b]{0.235\linewidth}\centering
    #2\\[2pt]{\small\texttt{#1}}\par\medskip
  \end{minipage}\hfill}

% the drawings of the manual
\setchemfig{atom sep=2em}

\title{\sffamily\bfseries\color{cfaccent}chemfrag\\[0.3em]
  \large\color{black}\mdseries Ready-made rings, chains and functional groups for \pkg{chemfig}}
\author{Vaibhav Blayer}
\date{Version \chemfragversion\ --- 2026-11-25}

\begin{document}
\maketitle
\thispagestyle{empty}

\begin{center}
\setchemfig{atom sep=1.8em}
\chemfig{!{benzene*}{2=NO_2, 4=Cl}-[::0]!{OH}}\hfill
\chemfig{H_3C-!{benzene*}{exit=4}-[::0]!{carboxyl}}\hfill
\chemfig{!{chain}{6, 2=CH_3, 4=OH}}\hfill
\chemfig{!{pyridine*}{entry=3, circle}-[::0]!{COOH}}
\end{center}

\begin{abstract}
\pkg{chemfrag} adds building blocks to \pkg{chemfig}: rings, zigzag chains
and functional groups. A fragment is not a picture but a \pkg{chemfig}
\emph{submolecule} that turns into ordinary \pkg{chemfig} code. You keep
writing \cs{chemfig}\{\dots\} exactly as before and drop fragments in wherever
they help. You can continue with hand-written \pkg{chemfig} code from any
fragment, and attach fragments to any hand-written structure.
\end{abstract}

\newpage
\tableofcontents
\newpage

% =====================================================================
\section{Introduction}

Drawing a substituted benzene with \pkg{chemfig} means writing the ring by
hand and placing every substituent inside it:
\begin{example}
\chemfig{*6((-OH)=(-NO_2)-=(-Cl)-=-)}
\end{example}
That is precise, but you have to remember which bond comes first, where the
branches go and which way the ring is drawn. With \pkg{chemfrag} the same
molecule is a ring with options:
\begin{example}
\chemfig{!{benzene*}{1=OH, 2=NO_2, 4=Cl}}
\end{example}
The second line produces the \emph{same} \pkg{chemfig} code as the first.
That is the whole idea of the package:
\begin{itemize}
  \item \textbf{It never draws anything itself.} Every fragment expands to
        plain \pkg{chemfig} code. All \pkg{chemfig} settings (\cs{setchemfig},
        bond styles, \cs{chemmove}, schemes) keep working.
  \item \textbf{It mixes with hand-written code in both directions.} Start
        from a fragment and continue in \pkg{chemfig}, or start in
        \pkg{chemfig} and attach a fragment.
  \item \textbf{Fragments are numbered like in IUPAC names}, so
        ``2-nitro-4-chlorophenol'' turns into keys almost word for word.
\end{itemize}

\subsection{Installation}
\pkg{chemfrag} needs \pkg{chemfig} 1.6 or later and a recent \LaTeX{}
(2022 or later). Put the files \texttt{chemfrag.sty} and
\texttt{chemfrag-*.code.tex} where \TeX{} can find them (or install the
package from CTAN with your distribution), then load it:
\begin{code}
\usepackage{chemfrag}   % loads chemfig
\end{code}

% =====================================================================
\section{The five-minute tour}

\paragraph{Plain fragments.} \frag{benzene}, \frag{cyclohexane},
\frag{pyridine}, \frag{OH}, \frag{COOH}\dots{} are \pkg{chemfig}
submolecules. Use them anywhere in \cs{chemfig}:
\begin{example}
\chemfig{!{benzene}}\quad
\chemfig{!{cyclohexane}}\quad
\chemfig{!{furan}}
\end{example}

\paragraph{Fragments with options.} Add a star and a list of keys. Numbers
are positions on the ring:
\begin{example}
\chemfig{!{benzene*}{1=CH_3, 2=NO_2, 4=NO_2, 6=NO_2}}
\end{example}

\paragraph{Continuing with chemfig.} After a ring, \pkg{chemfig} goes on
from the \emph{exit atom} (atom~1 by default) and the bond
\verb|-[::0]| points straight out of the ring. From there you write normal
\pkg{chemfig}:
\begin{example}
\chemfig{!{benzene}-[::0]C(=[::60]O)-[::-60]OH}
\end{example}
or attach another fragment:
\begin{example}
\chemfig{!{benzene}-[::0]!{carboxyl}}
\end{example}

\paragraph{Joining fragments.} A ring reached through a bond leaves on the
opposite atom, so fragments line up:
\begin{example}
\chemfig{!{NO2}-!{benzene}-[::0]!{benzene}-[::0]!{NH2}}
\end{example}

\paragraph{Chains.} \frag{chain}\verb|{n}| draws a zigzag chain of $n$
atoms, with substituents and double bonds by number:
\begin{example}
\chemfig{!{chain}{5, 2=CH_3, double=3}-[::0]!{OH}}
\end{example}

\paragraph{Functional groups} read the right way round by themselves:
\begin{example}
\chemfig{!{COOH}-[:180]!{benzene*}{exit=4}-[::0]!{COOH}}
\end{example}

% =====================================================================
\section{Concepts}

\subsection{Plain and starred fragments}
Each ring \meta{name} exists twice:
\begin{center}
\begin{tabular}{>{\ttfamily}l l}
\toprule
!\{\meta{name}\}            & the ring with default settings\\
!\{\meta{name}*\}\{\meta{keys}\} & the ring with options \\
\bottomrule
\end{tabular}
\end{center}
The starred form \emph{always} needs its braces, even if they are empty.
\pkg{chemfig} cannot tell whether an argument follows a submolecule when
the submolecule ends the molecule, so the star tells it.
\frag{chain} always takes an argument and has no star.

\subsection{Numbering}
\begin{itemize}
  \item Atom \key{1} is the atom where the ring is entered: the atom a
        bond arrives at, or the first atom of a ring that starts the molecule.
  \item Numbers increase counter-clockwise. Use \key{clockwise} to reverse
        them.
  \item \key{entry=}\meta{n} gives the entry atom another number. This is
        how you attach a pyridine by its carbon~3, for example.
\end{itemize}
Every key that takes a position (substituents, \key{atoms}, \key{bonds},
\key{exit}, \key{name}) uses this numbering.

\begin{example}
\chemfig{H_3C-!{pyridine*}{entry=3}}\quad
\chemfig{!{benzene*}{clockwise, 2=OH, 3=Cl}}
\end{example}

\subsection{Where drawing continues: the exit atom}
\pkg{chemfig} itself always continues a ring from the atom where the ring
started. \pkg{chemfrag} adds invisible bonds after the ring: one to its
centre and one out to the \emph{exit atom}, or a single tiny one when the exit
atom is atom~1. The next bond therefore starts on the exit atom, and the
\emph{relative} angle \verb|[::0]| points straight out of the ring.
\begin{itemize}
  \item A ring that \textbf{starts the molecule} exits on atom~1.
  \item A ring \textbf{reached through a bond} exits on the opposite atom
        (atom~4 of a six-membered ring).
  \item \key{exit=}\meta{n} chooses any atom; \key{exit=none} removes the
        invisible bonds and gives exactly the ring you would write by hand.
        Use it for rings that nothing continues from: it is also faster.
\end{itemize}
\begin{example}
\chemfig{!{benzene*}{exit=2}-[::0]CH_3}\quad
\chemfig{!{benzene*}{exit=3}-[::0]CH_3}\quad
\chemfig{!{benzene*}{exit=4}-[::0]CH_3}
\end{example}

\begin{tcolorbox}[cfbox,colback=yellow!6,colframe=orange!40,title={\sffamily Why
\texttt{-[::0]} and not just \texttt{-}?},coltitle=black,colbacktitle=orange!15]
In \pkg{chemfig} a bond without an angle is horizontal (\verb|[:0]|), whatever
came before. That is useful in a straight chain but would cut through a
ring. The relative angle \verb|[::0]| means ``continue in the direction of
the last bond'', which after a fragment is straight outwards. You can also
use \verb|[::60]|, \verb|[::-60]| or any absolute angle such as
\verb|[:90]| or \verb|[2]|.
\end{tcolorbox}

\subsection{Substituents}
A key \meta{n}\verb|=|\meta{chemfig code} puts a branch on atom \meta{n}.
\begin{itemize}
  \item Without a bond of its own, the substituent gets a single bond, and
        \pkg{chemfig} points it away from the ring.
  \item It may start with its own bond: \verb|1={=O}|, \verb|2={~N}|,
        \verb|3={-[::30]R}|. Values that contain \verb|=| or \verb|,| must be
        put in braces.
  \item It can be any \pkg{chemfig} code, including other fragments.
\end{itemize}
\begin{example}
\chemfig{!{cyclohexane*}{1={=O}, 3=CH_3}}\quad
\chemfig{!{benzene*}{1=!{carboxyl}, 2={-O-[::60]!{acetyl}}}}
\end{example}

% =====================================================================
\section{Rings}

\subsection{Keys}
\begin{longtable}{>{\raggedright}p{0.27\linewidth} p{0.65\linewidth}}
\toprule
\textbf{Key} & \textbf{Meaning}\\\midrule\endhead
\key{\meta{n}=\meta{code}} & substituent on atom \meta{n} (see above)\\
\key{exit=\meta{n}|none} & atom where drawing continues; default 1 at the start
  of a molecule, otherwise the opposite atom\\
\key{entry=\meta{n}} & number of the atom where the ring is entered (default 1)\\
\key{clockwise} & number the atoms clockwise\\
\key{circle} & draw an aromatic circle instead of alternating double bonds\\
\key{atoms=\{\meta{n}=\meta{atom},\dots\}} & atoms of the ring, e.g.\
  \verb|atoms={1=N}|; replaces the built-in list\\
\key{bonds=\meta{pattern}} & bonds of the ring, one character per bond;
  bond \meta{i} joins atoms \meta{i} and \meta{i}+1, e.g.\ \verb|bonds={=-=---}|\\
\key{size=\meta{n}} & number of atoms (for \frag{ring*})\\
\key{name=\meta{prefix}} & name the atoms \meta{prefix}1, \meta{prefix}2,
  \dots{} for \cs{chemmove}\\
\bottomrule
\end{longtable}
Keys given to the fragment override \cs{chemfragset} defaults, which override
the built-in settings of the ring. A flag can be switched off again with
\key{=false}, e.g.\ \verb|circle=false|.

\subsection{Built-in rings}
The small numbers show the default numbering. Each drawing was made with
the starred fragment of the ring, with the numbers given as substituents.

\medskip
\noindent
% chemfrag-gallery.tex -- ring gallery for chemfrag-doc.tex (generated list)
\ringcell{benzene}{\chemfig{!{benzene*}{1={-[,0.42,,,draw=none]\ringnum{1}}, 2={-[,0.42,,,draw=none]\ringnum{2}}, 3={-[,0.42,,,draw=none]\ringnum{3}}, 4={-[,0.42,,,draw=none]\ringnum{4}}, 5={-[,0.42,,,draw=none]\ringnum{5}}, 6={-[,0.42,,,draw=none]\ringnum{6}}}}}
\ringcell{cyclopropane}{\chemfig{!{cyclopropane*}{1={-[,0.42,,,draw=none]\ringnum{1}}, 2={-[,0.42,,,draw=none]\ringnum{2}}, 3={-[,0.42,,,draw=none]\ringnum{3}}}}}
\ringcell{cyclobutane}{\chemfig{!{cyclobutane*}{1={-[,0.42,,,draw=none]\ringnum{1}}, 2={-[,0.42,,,draw=none]\ringnum{2}}, 3={-[,0.42,,,draw=none]\ringnum{3}}, 4={-[,0.42,,,draw=none]\ringnum{4}}}}}
\ringcell{cyclopentane}{\chemfig{!{cyclopentane*}{1={-[,0.42,,,draw=none]\ringnum{1}}, 2={-[,0.42,,,draw=none]\ringnum{2}}, 3={-[,0.42,,,draw=none]\ringnum{3}}, 4={-[,0.42,,,draw=none]\ringnum{4}}, 5={-[,0.42,,,draw=none]\ringnum{5}}}}}
\ringcell{cyclohexane}{\chemfig{!{cyclohexane*}{1={-[,0.42,,,draw=none]\ringnum{1}}, 2={-[,0.42,,,draw=none]\ringnum{2}}, 3={-[,0.42,,,draw=none]\ringnum{3}}, 4={-[,0.42,,,draw=none]\ringnum{4}}, 5={-[,0.42,,,draw=none]\ringnum{5}}, 6={-[,0.42,,,draw=none]\ringnum{6}}}}}
\ringcell{cycloheptane}{\chemfig{!{cycloheptane*}{1={-[,0.42,,,draw=none]\ringnum{1}}, 2={-[,0.42,,,draw=none]\ringnum{2}}, 3={-[,0.42,,,draw=none]\ringnum{3}}, 4={-[,0.42,,,draw=none]\ringnum{4}}, 5={-[,0.42,,,draw=none]\ringnum{5}}, 6={-[,0.42,,,draw=none]\ringnum{6}}, 7={-[,0.42,,,draw=none]\ringnum{7}}}}}
\ringcell{cyclooctane}{\chemfig{!{cyclooctane*}{1={-[,0.42,,,draw=none]\ringnum{1}}, 2={-[,0.42,,,draw=none]\ringnum{2}}, 3={-[,0.42,,,draw=none]\ringnum{3}}, 4={-[,0.42,,,draw=none]\ringnum{4}}, 5={-[,0.42,,,draw=none]\ringnum{5}}, 6={-[,0.42,,,draw=none]\ringnum{6}}, 7={-[,0.42,,,draw=none]\ringnum{7}}, 8={-[,0.42,,,draw=none]\ringnum{8}}}}}
\ringcell{cyclopentene}{\chemfig{!{cyclopentene*}{1={-[,0.42,,,draw=none]\ringnum{1}}, 2={-[,0.42,,,draw=none]\ringnum{2}}, 3={-[,0.42,,,draw=none]\ringnum{3}}, 4={-[,0.42,,,draw=none]\ringnum{4}}, 5={-[,0.42,,,draw=none]\ringnum{5}}}}}
\ringcell{cyclohexene}{\chemfig{!{cyclohexene*}{1={-[,0.42,,,draw=none]\ringnum{1}}, 2={-[,0.42,,,draw=none]\ringnum{2}}, 3={-[,0.42,,,draw=none]\ringnum{3}}, 4={-[,0.42,,,draw=none]\ringnum{4}}, 5={-[,0.42,,,draw=none]\ringnum{5}}, 6={-[,0.42,,,draw=none]\ringnum{6}}}}}
\ringcell{cyclopentadiene}{\chemfig{!{cyclopentadiene*}{1={-[,0.42,,,draw=none]\ringnum{1}}, 2={-[,0.42,,,draw=none]\ringnum{2}}, 3={-[,0.42,,,draw=none]\ringnum{3}}, 4={-[,0.42,,,draw=none]\ringnum{4}}, 5={-[,0.42,,,draw=none]\ringnum{5}}}}}
\ringcell{cyclohexadiene}{\chemfig{!{cyclohexadiene*}{1={-[,0.42,,,draw=none]\ringnum{1}}, 2={-[,0.42,,,draw=none]\ringnum{2}}, 3={-[,0.42,,,draw=none]\ringnum{3}}, 4={-[,0.42,,,draw=none]\ringnum{4}}, 5={-[,0.42,,,draw=none]\ringnum{5}}, 6={-[,0.42,,,draw=none]\ringnum{6}}}}}
\ringcell{pyridine}{\chemfig{!{pyridine*}{1={-[,0.42,,,draw=none]\ringnum{1}}, 2={-[,0.42,,,draw=none]\ringnum{2}}, 3={-[,0.42,,,draw=none]\ringnum{3}}, 4={-[,0.42,,,draw=none]\ringnum{4}}, 5={-[,0.42,,,draw=none]\ringnum{5}}, 6={-[,0.42,,,draw=none]\ringnum{6}}}}}
\ringcell{pyrimidine}{\chemfig{!{pyrimidine*}{1={-[,0.42,,,draw=none]\ringnum{1}}, 2={-[,0.42,,,draw=none]\ringnum{2}}, 3={-[,0.42,,,draw=none]\ringnum{3}}, 4={-[,0.42,,,draw=none]\ringnum{4}}, 5={-[,0.42,,,draw=none]\ringnum{5}}, 6={-[,0.42,,,draw=none]\ringnum{6}}}}}
\ringcell{pyrazine}{\chemfig{!{pyrazine*}{1={-[,0.42,,,draw=none]\ringnum{1}}, 2={-[,0.42,,,draw=none]\ringnum{2}}, 3={-[,0.42,,,draw=none]\ringnum{3}}, 4={-[,0.42,,,draw=none]\ringnum{4}}, 5={-[,0.42,,,draw=none]\ringnum{5}}, 6={-[,0.42,,,draw=none]\ringnum{6}}}}}
\ringcell{pyrrole}{\chemfig{!{pyrrole*}{1={-[,0.42,,,draw=none]\ringnum{1}}, 2={-[,0.42,,,draw=none]\ringnum{2}}, 3={-[,0.42,,,draw=none]\ringnum{3}}, 4={-[,0.42,,,draw=none]\ringnum{4}}, 5={-[,0.42,,,draw=none]\ringnum{5}}}}}
\ringcell{furan}{\chemfig{!{furan*}{1={-[,0.42,,,draw=none]\ringnum{1}}, 2={-[,0.42,,,draw=none]\ringnum{2}}, 3={-[,0.42,,,draw=none]\ringnum{3}}, 4={-[,0.42,,,draw=none]\ringnum{4}}, 5={-[,0.42,,,draw=none]\ringnum{5}}}}}
\ringcell{thiophene}{\chemfig{!{thiophene*}{1={-[,0.42,,,draw=none]\ringnum{1}}, 2={-[,0.42,,,draw=none]\ringnum{2}}, 3={-[,0.42,,,draw=none]\ringnum{3}}, 4={-[,0.42,,,draw=none]\ringnum{4}}, 5={-[,0.42,,,draw=none]\ringnum{5}}}}}
\ringcell{imidazole}{\chemfig{!{imidazole*}{1={-[,0.42,,,draw=none]\ringnum{1}}, 2={-[,0.42,,,draw=none]\ringnum{2}}, 3={-[,0.42,,,draw=none]\ringnum{3}}, 4={-[,0.42,,,draw=none]\ringnum{4}}, 5={-[,0.42,,,draw=none]\ringnum{5}}}}}
\ringcell{oxazole}{\chemfig{!{oxazole*}{1={-[,0.42,,,draw=none]\ringnum{1}}, 2={-[,0.42,,,draw=none]\ringnum{2}}, 3={-[,0.42,,,draw=none]\ringnum{3}}, 4={-[,0.42,,,draw=none]\ringnum{4}}, 5={-[,0.42,,,draw=none]\ringnum{5}}}}}
\ringcell{thiazole}{\chemfig{!{thiazole*}{1={-[,0.42,,,draw=none]\ringnum{1}}, 2={-[,0.42,,,draw=none]\ringnum{2}}, 3={-[,0.42,,,draw=none]\ringnum{3}}, 4={-[,0.42,,,draw=none]\ringnum{4}}, 5={-[,0.42,,,draw=none]\ringnum{5}}}}}
\ringcell{oxirane}{\chemfig{!{oxirane*}{1={-[,0.42,,,draw=none]\ringnum{1}}, 2={-[,0.42,,,draw=none]\ringnum{2}}, 3={-[,0.42,,,draw=none]\ringnum{3}}}}}
\ringcell{tetrahydrofuran}{\chemfig{!{tetrahydrofuran*}{1={-[,0.42,,,draw=none]\ringnum{1}}, 2={-[,0.42,,,draw=none]\ringnum{2}}, 3={-[,0.42,,,draw=none]\ringnum{3}}, 4={-[,0.42,,,draw=none]\ringnum{4}}, 5={-[,0.42,,,draw=none]\ringnum{5}}}}}
\ringcell{pyrrolidine}{\chemfig{!{pyrrolidine*}{1={-[,0.42,,,draw=none]\ringnum{1}}, 2={-[,0.42,,,draw=none]\ringnum{2}}, 3={-[,0.42,,,draw=none]\ringnum{3}}, 4={-[,0.42,,,draw=none]\ringnum{4}}, 5={-[,0.42,,,draw=none]\ringnum{5}}}}}
\ringcell{piperidine}{\chemfig{!{piperidine*}{1={-[,0.42,,,draw=none]\ringnum{1}}, 2={-[,0.42,,,draw=none]\ringnum{2}}, 3={-[,0.42,,,draw=none]\ringnum{3}}, 4={-[,0.42,,,draw=none]\ringnum{4}}, 5={-[,0.42,,,draw=none]\ringnum{5}}, 6={-[,0.42,,,draw=none]\ringnum{6}}}}}
\ringcell{piperazine}{\chemfig{!{piperazine*}{1={-[,0.42,,,draw=none]\ringnum{1}}, 2={-[,0.42,,,draw=none]\ringnum{2}}, 3={-[,0.42,,,draw=none]\ringnum{3}}, 4={-[,0.42,,,draw=none]\ringnum{4}}, 5={-[,0.42,,,draw=none]\ringnum{5}}, 6={-[,0.42,,,draw=none]\ringnum{6}}}}}
\ringcell{morpholine}{\chemfig{!{morpholine*}{1={-[,0.42,,,draw=none]\ringnum{1}}, 2={-[,0.42,,,draw=none]\ringnum{2}}, 3={-[,0.42,,,draw=none]\ringnum{3}}, 4={-[,0.42,,,draw=none]\ringnum{4}}, 5={-[,0.42,,,draw=none]\ringnum{5}}, 6={-[,0.42,,,draw=none]\ringnum{6}}}}}
\ringcell{oxane}{\chemfig{!{oxane*}{1={-[,0.42,,,draw=none]\ringnum{1}}, 2={-[,0.42,,,draw=none]\ringnum{2}}, 3={-[,0.42,,,draw=none]\ringnum{3}}, 4={-[,0.42,,,draw=none]\ringnum{4}}, 5={-[,0.42,,,draw=none]\ringnum{5}}, 6={-[,0.42,,,draw=none]\ringnum{6}}}}}


\medskip
Fused rings are available as plain fragments without options:
\frag{naphthalene}, \frag{anthracene} and \frag{indole}.
\begin{example}
\chemfig{!{naphthalene}}\quad
\chemfig{!{indole}}
\end{example}

\subsection{Any ring}
\frag{ring*} draws any ring from scratch:
\begin{example}
\chemfig{!{ring*}{size=7, bonds={=-=-=--}, 3=Br}}
\end{example}

\subsection{More examples}
\begin{example}
\chemfig{!{benzene*}{circle, 1=OH, 4=NO_2}}
\end{example}
\begin{example}
\chemfig{!{furan*}{2=CHO}}\quad
\chemfig{!{thiophene*}{2=CH_3}}
\end{example}
\begin{example}
\chemfig{!{piperidine*}{1=, exit=1}-[::0]!{acetyl}}
\end{example}

% =====================================================================
\section{Chains}

\frag{chain}\verb|{|\meta{keys}\verb|}| draws a zigzag chain. It starts on the
current atom (C1), so \verb|R-!{chain}{3}| is a propyl group on R. All angles are
relative, so the chain follows whatever bond leads into it.

\begin{longtable}{>{\raggedright}p{0.27\linewidth} p{0.65\linewidth}}
\toprule
\textbf{Key} & \textbf{Meaning}\\\midrule\endhead
\key{\meta{n}} or \key{n=\meta{n}} & number of atoms (default 3)\\
\key{\meta{k}=\meta{code}} & substituent on atom \meta{k}\\
\key{start=up|down} & direction of the first zig (default \texttt{up})\\
\key{double=\{\dots\}}, \key{triple=\{\dots\}} & numbers of double / triple
  bonds; bond \meta{j} joins atoms \meta{j} and \meta{j}+1\\
\key{bonds=\meta{pattern}} & bond characters one by one\\
\key{atoms=\{\meta{k}=\meta{atom},\dots\}} & heteroatoms in the chain\\
\key{name=\meta{prefix}} & atom names for \cs{chemmove}\\
\key{exit=none} & no final turn (see below)\\
\bottomrule
\end{longtable}

After the last atom the chain turns invisibly towards the next zig, so
\verb|-[::0]| continues the zigzag.
\begin{example}
\chemfig{!{chain}{6, 2=CH_3, 4=OH}}
\end{example}
\begin{example}
\chemfig{!{chain}{4, double=2, start=down}}\quad
\chemfig{!{chain}{5, atoms={3=O}}}
\end{example}
\begin{example}
\chemfig{!{chain}{3, triple=2}-[::0]!{OH}}
\end{example}
\begin{example}
\chemfig{!{benzene}-[::0]!{chain}{3}-[::0]!{OH}}
\end{example}

% =====================================================================
\section{Functional groups}

\subsection{Condensed groups}
These groups are written as formulas. Like all \pkg{chemfig} submolecules
defined with two versions, they read from the ring outwards: when the bond
arrives from the right, the mirrored formula is used.

\medskip
\begin{center}\small
\begin{tabular}{>{\ttfamily}l c c @{\qquad} >{\ttfamily}l c c}
\toprule
\normalfont fragment & right & left & \normalfont fragment & right & left\\\midrule
!\{OH\}    & \chemfig{-!{OH}}    & \chemfig{-[4]!{OH}}    & !\{CN\}    & \chemfig{-!{CN}}    & \chemfig{-[4]!{CN}}\\
!\{SH\}    & \chemfig{-!{SH}}    & \chemfig{-[4]!{SH}}    & !\{COOH\}  & \chemfig{-!{COOH}}  & \chemfig{-[4]!{COOH}}\\
!\{NH2\}   & \chemfig{-!{NH2}}   & \chemfig{-[4]!{NH2}}   & !\{COOMe\} & \chemfig{-!{COOMe}} & \chemfig{-[4]!{COOMe}}\\
!\{NHMe\}  & \chemfig{-!{NHMe}}  & \chemfig{-[4]!{NHMe}}  & !\{COOEt\} & \chemfig{-!{COOEt}} & \chemfig{-[4]!{COOEt}}\\
!\{NMe2\}  & \chemfig{-!{NMe2}}  & \chemfig{-[4]!{NMe2}}  & !\{COCl\}  & \chemfig{-!{COCl}}  & \chemfig{-[4]!{COCl}}\\
!\{NO2\}   & \chemfig{-!{NO2}}   & \chemfig{-[4]!{NO2}}   & !\{CONH2\} & \chemfig{-!{CONH2}} & \chemfig{-[4]!{CONH2}}\\
!\{CHO\}   & \chemfig{-!{CHO}}   & \chemfig{-[4]!{CHO}}   & !\{COMe\}  & \chemfig{-!{COMe}}  & \chemfig{-[4]!{COMe}}\\
!\{OMe\}   & \chemfig{-!{OMe}}   & \chemfig{-[4]!{OMe}}   & !\{Me\}    & \chemfig{-!{Me}}    & \chemfig{-[4]!{Me}}\\
!\{OEt\}   & \chemfig{-!{OEt}}   & \chemfig{-[4]!{OEt}}   & !\{Et\}    & \chemfig{-!{Et}}    & \chemfig{-[4]!{Et}}\\
!\{OAc\}   & \chemfig{-!{OAc}}   & \chemfig{-[4]!{OAc}}   & !\{CF3\}   & \chemfig{-!{CF3}}   & \chemfig{-[4]!{CF3}}\\
!\{SO3H\}  & \chemfig{-!{SO3H}}  & \chemfig{-[4]!{SO3H}}  & !\{CCl3\}  & \chemfig{-!{CCl3}}  & \chemfig{-[4]!{CCl3}}\\
\bottomrule
\end{tabular}
\end{center}

\paragraph{At the start of a molecule} no bond arrives, so \pkg{chemfig}
uses the right-hand form. If the molecule goes off to the right, write
the formula yourself (\verb|HO-...|), or start the molecule with the bond
pointing left as in the tour (\verb|!{COOH}-[:180]...|).

\subsection{Structural groups}
These are drawn as skeletal formulas with relative angles, so they fit any
direction. The C=O goes to the left of the direction of travel. Add a prime
(\verb|'|) to the name for the mirror image.

\medskip
\begin{center}\small
\renewcommand\arraystretch{3.4}
\begin{tabular}{>{\ttfamily}l c @{\qquad} >{\ttfamily}l c}
\toprule
\normalfont fragment & drawing & \normalfont fragment & drawing\\\midrule
!\{carboxyl\}     & \chemfig{-!{carboxyl}}     & !\{carboxyl'\}     & \chemfig{-!{carboxyl'}}\\
!\{aldehyde\}     & \chemfig{-!{aldehyde}}     & !\{ester\}         & \chemfig{-!{ester}}\\
!\{amide\}        & \chemfig{-!{amide}}        & !\{acylchloride\}  & \chemfig{-!{acylchloride}}\\
!\{acetyl\}       & \chemfig{-!{acetyl}}       & !\{carbonyl\}      & \chemfig{-!{carbonyl}}\\
!\{nitro\}        & \chemfig{-!{nitro}}        & !\{nitrile\}       & \chemfig{-!{nitrile}}\\
!\{sulfonyl\}     & \chemfig{-!{sulfonyl}-}    &                    & \\
\bottomrule
\end{tabular}
\end{center}
All structural groups except \texttt{nitrile} and \texttt{sulfonyl} have a
primed version.

\begin{example}
\chemfig{!{benzene*}{exit=3}-[::0]!{ester'}}
\end{example}

% =====================================================================
\section{Customising}

\subsection{Defaults}
\cs{chemfragset}\verb|{|\meta{keys}\verb|}| sets default keys for every following
fragment in the current group. Keys given to a fragment still win.
\begin{example}
\chemfragset{circle}
\chemfig{!{benzene}}\quad
\chemfig{!{pyridine}}\quad
\chemfig{!{benzene*}{circle=false}}
\end{example}

\subsection{New rings}
\cs{chemfragdefinering}\verb|{|\meta{name}\verb|}{|\meta{keys}\verb|}| defines
\frag{\meta{name}} and \frag{\meta{name}*}. The keys are the defaults of the new
ring.
\begin{example}
\chemfragdefinering{oxetane}
  {size=4, bonds={----}, atoms={1=O}}
\chemfig{!{oxetane*}{3=CH_3}}
\end{example}

\subsection{New groups}
Groups are plain \pkg{chemfig}; define your own with \cs{definesubmol}. The
optional argument is the code used when the bond comes from the right:
\begin{example}
\definesubmol{NHAc}[AcHN]{NHAc}
\chemfig{!{benzene*}{exit=4}-[::0]!{NHAc}}
\end{example}

\subsection{The engine behind the fragments}
The fragments are thin wrappers around two macros:
\begin{center}
\begin{tabular}{>{\ttfamily}l@{\quad is short for\quad}>{\ttfamily}l}
!\{benzene*\}\{...\} & !\textbackslash chemfragring\{benzene\}\{...\}\\
!\{chain\}\{...\}    & !\textbackslash chemfragchain\{...\}
\end{tabular}
\end{center}
Both macros are fully expandable and expand to \pkg{chemfig} code:
\begin{code}
\chemfragring{benzene}{2=NO_2, 4=Cl}
% gives (in the middle of a molecule)
*6(=(-NO_2)-=(-Cl)-=-)-[::0,1,,,draw=none]-[::0,1,,,draw=none]
\end{code}

% =====================================================================
\section{Reaction mechanisms}
\key{name=}\meta{prefix} names the atoms of a ring or chain, so you can
draw arrows between them with \cs{chemmove} (two compilations needed):
\begin{example}
\chemfig{!{benzene*}{name=r, 1=OMe}}
\chemmove[->,red,shorten <=2pt,shorten >=2pt]{
  \draw (r2) .. controls +(-90:9mm)
        and +(-30:9mm) .. (r3);}
\end{example}

% =====================================================================
\section{Worked examples}

\subsection*{Aspirin}
\begin{example}
\chemfig{!{benzene*}{2={-O-[::60]!{acetyl'}}}
  -[::0]!{carboxyl}}
\end{example}

\subsection*{Paracetamol}
\begin{example}
\chemfig{HO-!{benzene}-[::0]N(-[::90]H)
  -[::-60]!{acetyl'}}
\end{example}

\subsection*{Ibuprofen}
\begin{examplewide}
\chemfig{!{chain}{3, 2=CH_3}-[::0]!{benzene}
  -[::0](-[::90]CH_3)-[::-60]!{carboxyl}}
\end{examplewide}

\subsection*{Vanillin}
\begin{example}
\chemfig{!{benzene*}{2=OMe, 4=CHO}-[::0]!{OH}}
\end{example}

\subsection*{Nicotinic acid}
\begin{example}
\chemfig{!{pyridine*}{entry=3}-[::0]!{carboxyl}}
\end{example}

\subsection*{A reaction}
\begin{examplewide}
\schemestart
  \chemfig{!{benzene}-[::0]!{OH}}
  \arrow{->[\chemfig{HNO_3}]}
  \chemfig{!{benzene*}{4=NO_2}-[::0]!{OH}}
  \+
  \chemfig{!{benzene*}{2=NO_2}-[::0]!{OH}}
\schemestop
\end{examplewide}

% =====================================================================
\section{Notes and limitations}
\begin{itemize}
  \item Fragments work with \pkg{babel}'s French and other languages that
        make \verb|!| or \verb|:| active, as long as \cs{chemfig} itself works.
  \item Values with \verb|=| or \verb|,| must be braced:
        \verb|1={=O}|, \verb|atoms={1=N,3=N}|.
  \item A substituent on the entry atom of a ring that is reached through a
        bond is ignored, because the incoming bond takes its place.
  \item \pkg{chemfrag} reads \pkg{chemfig}'s internal state
        (\verb|\CF_lastaction|) to tell whether a ring starts the molecule. If
        a future \pkg{chemfig} removes it, rings behave as if a bond arrived
        and still draw correctly.
  \item \textbf{Speed.} A \pkg{chemfig} molecule is much slower to typeset
        than the same lines drawn directly in TikZ, because \pkg{chemfig}
        computes every bond (about 7--12\,ms each with pdf\LaTeX).
        \pkg{chemfrag}'s own expansion costs about 2\,ms per fragment. Each
        invisible helper bond costs as much as a real bond, so a ring with an
        exit is slower than the hand-written ring:
        one helper bond for an exit on atom~1, two for any other exit, plus a
        tiny entry bond when the ring starts the molecule. A chain adds one
        helper bond at its end. Write \key{exit=none} when nothing continues
        from a fragment. In documents with many molecules, typeset each one
        only once with the \pkg{memoize} package, TikZ externalization or the
        \pkg{standalone} class.
\end{itemize}

% =====================================================================
\section{Change history}
\begin{description}[font=\sffamily\bfseries]
  \item[1.0.1 (2026-11-25)] Fewer invisible helper bonds: an exit on atom~1
    uses one bond instead of two, and a ring with \key{exit=none} has none,
    so it is identical to the hand-written ring. Drawings are unchanged.
  \item[1.0.0 (2026-11-25)] First public release: 27 rings, generic rings,
    chains, 23 condensed and 18 structural groups, \cs{chemfragset},
    \cs{chemfragdefinering}, atom names for \cs{chemmove}.
\end{description}

\end{document}
